\section{Methodology}
\label{sec:method}

We follow the contrastive activation analysis pipeline established by \citet{arditi2024refusal} and \citet{marks2024geometry}, adapted for stylistic properties.
The approach has four stages: data construction, activation collection, direction finding, and validation.

\subsection{Data Construction}
\label{sec:data}

We use the \hcthree dataset~\citep{guo2023hc3}, which contains paired human and ChatGPT responses to the same questions.
The content-controlled pairing is critical: because both responses answer the same question, differences in activation patterns are more likely to reflect \emph{style} rather than \emph{topic}.

From \hcthree, we extract 392 paired samples after filtering for a minimum length of 10 words.
We select the first human and first ChatGPT answer for each question, truncate to 512 characters, and shuffle with a fixed seed.
The samples span five domains: \textsc{reddit\_eli5} (271 samples), \textsc{finance} (74), \textsc{medicine} (18), \textsc{open\_qa} (15), and \textsc{wiki\_csai} (14).
The domain imbalance is a known property of \hcthree and motivates our cross-domain generalization experiments.

\subsection{Activation Collection}
\label{sec:activations}

We process each text through the model using TransformerLens and collect residual stream activations at the last token position for every layer.
We use the last token because it aggregates information from the full sequence, following the standard protocol for sequence-level probing~\citep{marks2024geometry}.

We test two models:
\begin{itemize}[leftmargin=*,itemsep=0pt,topsep=0pt]
    \item \pythia: a 32-layer base model ($d_\text{model} = 2560$) with no alignment training, serving as a control for whether the signal exists prior to instruction tuning.
    \item \qwen: a 36-layer instruction-tuned model ($d_\text{model} = 2048$) that has undergone RLHF, representing a model that actively \emph{produces} \aistyle text.
\end{itemize}

All inference uses float16 precision with batch size 8 and a maximum of 128 tokens per input.

\subsection{Direction Finding}
\label{sec:directions}

\para{Difference-in-means (DiffMean).}
For each layer $l$, we compute the candidate \aistyle direction as:
\begin{equation}
    \vd_l = \frac{1}{|\gS_\text{AI}|} \sum_{i \in \gS_\text{AI}} \vh_l^{(i)} - \frac{1}{|\gS_\text{H}|} \sum_{i \in \gS_\text{H}} \vh_l^{(i)}
    \label{eq:diffmean}
\end{equation}
where $\vh_l^{(i)} \in \R^{d_\text{model}}$ is the residual stream activation at layer $l$ for sample $i$, and $\gS_\text{AI}$, $\gS_\text{H}$ are the sets of AI-written and human-written samples, respectively.
We choose \diffmean over logistic regression because \citet{marks2024geometry} showed it identifies more causally relevant directions.

\para{Linear probes.}
As a stronger classification baseline, we train an SGDClassifier (logistic regression with stochastic gradient descent, $\alpha = 10^{-4}$) on the full-dimensional residual stream activations.
We evaluate with 5-fold stratified cross-validation.

\subsection{Validation}
\label{sec:validation}

We validate the identified direction along four axes.

\para{Random direction baseline.}
We sample 500 random unit vectors in $\R^{d_\text{model}}$ and compute each one's classification accuracy when used as a 1D projection (threshold at median).
We compare the \diffmean direction's accuracy to this null distribution and report the $z$-score and $p$-value.

\para{Cross-domain generalization.}
We train the \diffmean direction on samples from one domain and test on each other domain.
A direction that captures general \aistyle rather than domain-specific artifacts should transfer across domains.

\para{Causal steering.}
We add a scaled version of the \diffmean direction to the residual stream during generation:
\begin{equation}
    \vh_l' = \vh_l + \alpha \cdot \hat{\vd}_l
    \label{eq:steering}
\end{equation}
where $\hat{\vd}_l$ is the unit-normalized direction and $\alpha \in \{-20, -10, -5, 0, +5, +10, +20\}$.
We measure the resulting projection of generated text activations onto the direction and observe qualitative changes in output style.

\para{Layer analysis.}
We compare probe accuracy and \diffmean accuracy across all layers to identify where the model encodes the \aistyle signal most strongly.
